\documentclass[../main.tex]{subfiles}

\begin{document}

\chapter{Fundamental Math for Machine Learning}

This note covers fundamental math necessary for machine learning.
Please note that I won't nor have the capacity to include details
on the concepts that are not strictly necessary to learn ML,
which I believe are out of the scope of the notes. With that
in mind, let's get started!

\section{Calculus}

Ah yes, that one friend who's everywhere, where ever you go.
Broadly speaking, calculus is one of the foundational math
for machine learning. For instance, backpropagation, which
we will discuss further in latter notes, is fundamentally
based on the chain rule. There are actually a lot to cover on
calculus and there are some advanced topics. Good news is that
you probably don't need calculus concepts outside multivariable
calculus (a.k.a. Calculus III) unless you are doing theoretical
works, and they are out of my capacity. That being said, I will
include most calculus concepts that you will need for general
purpose, but concepts that are out of scope of these notes will
not be included.

\subsection{Foundations}

I am assuming most of you have some background in calculus, so I will
be skipping basic concepts from Calculus AB/BC (a.k.a Calculus I and
II) as this section is about calculus for machine learning.
Please don't worry if you don't as you probably won't need all
concepts from the courses to understand the topics here. Now, let's
start by building from terminologies.

\begin{definition}{}{}
The \textbf{limit} of a function is the value that the function
approaches as the input gets arbitrarily close to a value.
\end{definition}

Below are some examples.
\[
\lim_{x \to 3} x^2 = 3^2 = 9,
\quad
\lim_{x \to \infty} \frac{1}{x} = 0,
\quad
\lim_{x \to 1} \frac{1}{x - 1} \text{ DNE}
\]

\begin{definition}{}{}
\textbf{Derivative} of a function respect to a variable is the
measure of the rate of change of the output based on the change of
the variable(s). \textbf{Gradient} is a vector of partial derivatives.
\textbf{Differentiation} is the process of finding the derivative.
\end{definition}

With equations, derivative and partial derivatives can be written
as the following for function $f(x)$ and $f(x_1, \ldots, x_n)$.
\begin{align*}
    f'(a)
    &= \lim_{h \to 0} \frac{f(a + h) - f(a)}{h} \\
    \frac{\partial}{\partial x_i} f(a)
    &= \lim_{h \to 0}
        \frac{f(a_1, \ldots, a_{i-1}, a_i + h, a_{i+1}, \ldots, a_n)
                - f(a_1, \ldots, a_n)}{h}
\end{align*}
For gradient $\nabla f$, read as ``nabla f'', the mathematical
expression follows.
\[
\nabla f(a) =
\begin{bmatrix}
    \frac{\partial f}{\partial x_1} (a) \\
    \vdots \\
    \frac{\partial f}{\partial x_n} (a)
\end{bmatrix}
\]
For instance, the gradient of $f(x, y) = x^2 + 2y$ can be written as
$\nabla f = \langle 2x, 2 \rangle$ and the gradient at $(1, 1)$ 
can be written as $\nabla f(1, 1) = \langle 2, 2 \rangle$.

Now that we discussed about derivatives, let's take a look at common
rules that will appear for machine learning. Starting with the
power rule, below is the list of common rules.

\begin{theorem}
{Common Rules for Derivatives}
{Common Rules for Derivatives}
\begin{enumerate}[noitemsep]
\item $\frac{d}{dx} x^n = n x^{n-1}$
\item $\frac{d}{dx} e^x = e^x$
\item $\frac{d}{dx} \ln(x) = \frac{1}{x}$
\item $\frac{d}{dx} a = 0$
\item $\frac{d}{dx} cf(x) = cf'(x)$
\item $\frac{d}{dx} \left( f(x) \pm g(x) \right)
       = \frac{d}{dx} f(x) \pm \frac{d}{dx} g(x)$
\item $\frac{d}{dx} \left( f(x) g(x) \right)
       = f'(x)g(x) + f(x)g'(x)$
\item $\frac{d}{dx} \left( \frac{f(x)}{g(x)} \right)
       = \frac{f'(x)g(x) - f(x)g'(x)}{\left(g(x)\right)^2}$
\item $\frac{\partial^2 f}{\partial x^2}
       = \frac{\partial}{\partial x}
            \left( \frac{\partial f}{\partial x} \right)$
\item $\frac{\partial^2 f}{\partial y \partial x}
       = \frac{\partial}{\partial y}
            \left( \frac{\partial f}{\partial x} \right)$
\end{enumerate}
\end{theorem}

Let's skip the proofs for this one! In addition to derivative rules,
one of the most primary use case of calculus in machine learning is
the chain rule. Indeed, the neural network is fundamentally based on
the chain rule in calculus.

\begin{theorem}
{The Chain Rule of Calculus}
{The Chain Rule of Calculus}
The derivative of a composite function $h(x) = f \circ g (x)$ can be
represented as $h'(x) = f'(g(x)) \cdot g'(x)$. Alternatively,
for $y = f(u)$ and $u = g(x)$, the following equation holds.
\[
\frac{dy}{dx} = \frac{dy}{du} \cdot \frac{du}{dx}
\]
\end{theorem}
\begin{proof}
Consider the following equation by the definition of derivatives.
\[
h'(x)
= \lim_{n \to 0} \frac{h(x + n) - h(x)}{n}
\]
Substituting $h(x) = f(g(x))$, the following equations are obtained.
\begin{align*}
h'(x)
&= \lim_{n \to 0} \frac{f(g(x + n)) - f(g(x))}{n} \\
&= \lim_{n \to 0} \frac{f(g(x + n)) - f(g(x))}{g(x + n) - g(x)} \cdot
    \lim_{n \to 0} \frac{g(x + n) - g(x)}{n}
\end{align*}
Let $n' = g(x + n) - g(x)$. Assuming that $g$ is continuous,
as $n \to 0$, $n' \to 0$. Therefore, the following equation is
obtained.
\[
h'(x)
= \lim_{n' \to 0} \frac{f(u + n') - f(u)}{n'} \cdot
    \lim_{n \to 0} \frac{g(x + n) - g(x)}{n}
= \frac{dy}{du} \cdot \frac{du}{dx}
\]
Thus, the chain rule holds.
\end{proof}

For a brief example, let's take a look at an example with sigmoid
function.

\begin{problem}
{The Derivative of Sigmoid Function}
{The Derivative of Sigmoid Function}
Sigmoid function is defined as $\sigma(x) = \frac{1}{1 + e^{-x}}$.
Find the derivative of $\sigma(x)$.
\end{problem}
\begin{solution}
Consider the following process with the product and chain rule.
\begin{align*}
    \sigma'(x)
    &= \frac{d}{dx} \left( 1 + e^{-x} \right)^{-1}
    = - \left( 1 + e^{-x} \right)^{-2} \cdot \left( -e^{-x} \right) \\
    &= \left( 1 + e^{-x} \right)^{-2} \cdot e^{-x}
    = \left( 1 + e^{-x} \right)^{-1} \cdot
        \frac{e^{-x}}{1 + e^{-x}} \\
    &= \sigma(x) \left( 1 - \sigma(x) \right)
\end{align*}
Therefore, the derivative $\sigma'(x) = \sigma(x) (1 - \sigma(x))$.
\end{solution}

In addition to the chain rule above, we could expand it for
multivariable functions.

\begin{theorem}
{The Chain Rule for One Independent Variable}
{The Chain Rule for One Independent Variable}
For differentiable functions $x = a(t)$, $y = b(t)$, and
$z = f(x, y)$, the following equation holds.
\[
\frac{dz}{dt}
= \frac{\partial z}{\partial x} \cdot \frac{dx}{dt} +
    \frac{\partial z}{\partial y} \cdot \frac{dy}{dt}
\]
\end{theorem}

Similarly, we could expand it for cases where $x$ and $y$ are
functions of two variables.

\begin{theorem}
{The Chain Rule for Two Independent Variable}
{The Chain Rule for Two Independent Variable}
For differentiable functions $x = a(u, v)$, $y = b(u, v)$,
and $z = f(x, y)$, the following equations hold.
\begin{align*}
    \frac{\partial z}{\partial u}
    &= \frac{\partial z}{\partial x} \frac{\partial x}{\partial u} +
        \frac{\partial z}{\partial y} \frac{\partial y}{\partial u} \\
    \frac{\partial z}{\partial v}
    &= \frac{\partial z}{\partial x} \frac{\partial x}{\partial v} +
        \frac{\partial z}{\partial y} \frac{\partial y}{\partial v}
\end{align*}
\end{theorem}

For the purpose of these notes, let's skip the proof for the two
theorems above.



critical points
Linaer approximation, taylor approx, tangent planes

Vector calc

Jacobian and Hessian Matrices
Lagrange Multipliers
Fourier Transformation



\end{document}


